\section*{الفصل \ref{ch:b1:earthquake-source} --- مصدر الزلزال}

\begin{solution}{exo:b1:earthquake-source:1}
(أ) $\log_{10}(\num{4e17}) = 17.60$، إذن $M_w = \tfrac23(17.60 - 9.1) = 5.67$.
(ب) $\log_{10}(\num{5.59e22}) = 22.75$ و$M_w = \tfrac23(22.75 - 9.1) =
9.10$. (ج) $M_0 = 10^{1.5\times7.0 + 9.1} = 10^{19.6} = \qty{3.98e19}{N.m}$.
\end{solution}

\begin{solution}{exo:b1:earthquake-source:2}
(أ) $T = D/\dot s = \qty{3}{m}/\qty{0.025}{m.yr^{-1}} = \qty{120}{yr}$.
(ب) مع $c = 0.5$، $T = D/(c\dot s) = \qty{240}{yr}$: الزحف يستوعب نصف الحركة، والزلازل، بالانزلاق نفسه، تأتي نصف مرة كما كانت.
\end{solution}

\begin{solution}{exo:b1:earthquake-source:3}
بحسب \cref{prop:b1:earthquake-source:moment-steps}، نسبة العزمين هي
$10^{1.5\times(7.4 - 6.2)} = 10^{1.8} = 63.1$: يلزم 63 زلزالًا قدره 6.2 لإطلاق عزم زلزال واحد قدره 7.4.
\end{solution}

\begin{solution}{exo:b1:earthquake-source:4}
مع $b = 1$، كل وحدة قدرية تقسم العدد على 10: 10 زلازل قدرها 5 أو أكثر في السنة، و\num{1} قدرها 6 أو أكثر، و\num{0.1} قدرها 7 أو أكثر، أي واحد كل 10 سنوات في المتوسط.
\end{solution}

\begin{solution}{exo:b1:earthquake-source:5}
(أ) $A = \num{4.3e5} \times \qty{1.5e4}{m} = \qty{6.45e9}{m^2}$، و
$M_0 = \num{3e10}\times\num{6.45e9}\times4 = \qty{7.74e20}{N.m}$؛
$\log_{10}M_0 = 20.89$ و$M_w = \tfrac23(20.89 - 9.1) = 7.86$. (ب) عمق قدره
\qty{10}{km} يقسم المساحة والعزم على 1.5، وهذا يخفّض القدر بمقدار
$\tfrac23\log_{10}1.5 = 0.12$، إلى 7.74. (ج) الأول، 7.86، يُقرَّب إلى 7.9 من الكتالوج؛ والثاني أدنى قليلًا. والقدر ليس حسّاسًا كثيرًا للعمق: فعامل 1.5 على المساحة لا يساوي إلا 0.12 من القدر.
\end{solution}

\begin{solution}{exo:b1:earthquake-source:6}
(أ) $D = M_0/(\mu A) = \num{5.59e22}/(\num{3e10}\times\num{4e5}\times\num{1.5e5})
= \qty{31.1}{m}$. (ب) $\qty{31.1}{m}/\qty{0.083}{m.yr^{-1}} = \qty{374}{yr}$.
(ج) مع كسر قفل قدره 0.8، $31.1/(0.8\times0.083) = \qty{468}{yr}$؛ ومع
0.5، \qty{749}{yr}. فكلما قلّ قفل الواجهة، طال الزمن الذي يتراكم فيه انزلاق زلزال كبير واحد؛ ومتوسط الانزلاق الموجود هنا بمستطيل \qtyproduct{400 x 150}{km} و$\mu = \qty{30}{GPa}$ ليس إلا تقديرًا، والانزلاق الحقيقي لم يكن موحّدًا البتّة.
\end{solution}

\begin{solution}{exo:b1:earthquake-source:7}
(أ) \omterm{def:b1:geological-maps:attitude}{الامتداد} هو $\phi = 329$، \omterm{def:b1:geological-maps:attitude}{واتجاه الانحدار} $a = 059$، \omterm{def:b1:geological-maps:attitude}{والانحدار} $\delta =
\ang{8}$. ثم $\mathbf s = (\sin329, \cos329, 0) = (-0.5150, 0.8572, 0)$،
$\mathbf u = (-\cos8\sin59, -\cos8\cos59, \sin8) = (-0.8488, -0.5100, 0.1392)$،
$\mathbf n = (\sin8\sin59, \sin8\cos59, \cos8) = (0.1193, 0.0717, 0.9903)$،
و مع $r = \ang{110}$،
\[
  \mathbf d = \cos110\,\mathbf s + \sin110\,\mathbf u = (-0.6215, -0.7724, 0.1308) .
\]
(ب) \omterm{def:b1:earthquake-source:focal-mechanism}{المستوى المساعد} له العمود $\mathbf d$ والانزلاق $\mathbf n$.
بأخذ العمود صاعدًا، يكون انحداره $\arccos 0.1308 = \ang{82.5}$ نحو
$\arctan(-0.6215/-0.7724) = 218.8$، فيكون امتداده $128.8$؛ والانزلاق
$\mathbf n$ مُسقطًا على متجهتي \omterm{def:b1:geological-maps:attitude}{الامتداد} \omterm{def:b1:geological-maps:attitude}{والانحدار} الصاعد يعطي \omterm{def:b1:faults-folds-fabrics:fault-elements}{زاوية الانزلاق}
\ang{87.2}. مُقرَّبًا: $129/83/87$، وهو المستوى الثاني في الحلّ المنشور.
(ج) $\mathbf P = (\mathbf n - \mathbf d)/\sqrt2$ اتجاهه $221$ وغوصه
\ang{37}؛ $\mathbf T = (\mathbf n + \mathbf d)/\sqrt2$ اتجاهه $036$ وغوصه
\ang{52} (المنشور: $222/38$ و$36/52$). (د) المستوى الضحل، الذي ينحدر
\ang{8} نحو $059$، على طول الصفيحة الغائصة: \omterm{def:b1:faults-folds-fabrics:fault-types}{صدع معكوس} بزاوية انزلاق
\ang{110}، وبانحدار أقلّ من \ang{45}، صدع دفعي، هو واجهة الصفيحتين في منطقة الاندساس.
\end{solution}

\begin{solution}{exo:b1:earthquake-source:8}
(أ) المستوى ينحدر نحو $135$، ف$\mathbf n = (\sin135, \cos135, 0) =
(0.707, -0.707, 0)$؛ وانزلاق \omterm{def:b1:faults-folds-fabrics:fault-types}{أيمن} له $r = \ang{180}$، و$\mathbf d = -\mathbf s
= (-0.707, -0.707, 0)$. نحو الشمال، $\mathbf x = (0, 1, 0)$:
$(\mathbf n\cdot\mathbf x)(\mathbf d\cdot\mathbf x) = (-0.707)(-0.707) > 0$، دفعة.
شرقًا: $(0.707)(-0.707) < 0$، سحب. جنوبًا: دفعة؛ غربًا: سحب.
(ب) $\mathbf P = (\mathbf n - \mathbf d)/\sqrt2 = (1, 0, 0)$: \omterm{def:b1:earthquake-source:focal-mechanism}{محور P} أفقي
ويتّجه $090$--$270$؛ و$\mathbf T = (0, -1, 0)$ يتّجه
$000$--$180$. (ج) الصدع الشاقولي الذي يمتدّ $135$--$315$ عموده $\mathbf d$. ومع الانزلاق $\mathbf n$، فالكتلة الواقعة على جانبه الجنوبي الغربي، التي يشير إليها $\mathbf d$، تتحرّك نحو $135$؛ ومن الكتلة الشمالية الشرقية، ناظرين نحو الجنوب الغربي، يكون ذلك إلى اليسار: انزلاق \omterm{def:b1:faults-folds-fabrics:fault-types}{أيسر}. وبحسب
\cref{prop:b1:earthquake-source:auxiliary}، يعطي هذا الانزلاق الحركات الأولى نفسها في كل اتجاه.
\end{solution}

\begin{solution}{exo:b1:earthquake-source:9}
(أ) $V = \qty{40}{mm.yr^{-1}}$ و$v = V/4$ عند \qty{10}{km}، فبحسب
\cref{prop:b1:earthquake-source:profile}(ب) $D = \qty{10}{km}$؛ مباشرة،
$D = x/\tan(\pi v/V) = 10/\tan(\pi/4) = \qty{10}{km}$. (ب) $v(30) =
(40/\pi)\arctan 3 = \qty{15.9}{mm.yr^{-1}}$. (ج) $V/(\pi D) =
\num{0.040}/(\pi\times\num{10000}) = \num{1.27e-6}$ في السنة.
\end{solution}

\begin{solution}{exo:b1:earthquake-source:10}
(أ) عزوم الفئات ما زالت $C q^k$ مع $q = 10^{1.5 - b}$، لكن الآن $q < 1$: عزم الفئات يتناقص نحو القدرات الكبيرة، والفئة الدنيا، $k = 0$، تحمل النسبة $C/S = (1 - q)/(1 - q^{K+1}) > 1 - q$، وهي الصورة المرآتية ل\cref{prop:b1:earthquake-source:largest}. فالزلازل الأصغر المحسوبة تهيمن عندئذٍ، ويكبر المجموع دون حدّ كلما أُدخلت زلازل أصغر فأصغر. (ب) $1 - 10^{-0.7} = 0.80$: على الأقلّ \qty{80}{\percent}. (ج) $\log_{10}N = 8.20 - 1.01\times9 = -0.89$، إذن $N = 0.13$
في السنة وقرابة 6.4 في خمسين سنة، ثلاثة أضعاف الزلزالين المرصودين. والخطّ، المطابَق بين القدرين 6 و\num{7.5}، يبالغ في تقدير أكبر الزلازل؛ إما أن القانون ينحني نحو الأسفل عند القمة، كما يوحي الحجم المحدود للصدوع، أو أن خمسين سنة أقلّ من أن تعاين أحداثًا تقع مرة كل عقود. والاثنان ممكنان مع زلزالين اثنين.
\end{solution}

\begin{solution}{exo:b1:earthquake-source:11}
(أ) بعد 150 سنة يتحرّك عمود عند $x$ بمقدار $150\,v(x) = 150\times(34/\pi)
\arctan(x/15)$ ملليمترًا بالنسبة للعمود الواقع على الصدع: \qty{0.11}{m} عند
\qty{1}{km}، \qty{1.28}{m} عند \qty{15}{km}، \qty{2.08}{m} عند \qty{50}{km}،
نحو الجنوب لصدع \omterm{def:b1:faults-folds-fabrics:fault-types}{أيمن} يُرى من الغرب. (ب) عجز الانزلاق هو
$s = \qty{0.034}{m.yr^{-1}}\times\qty{150}{yr} = \qty{5.1}{m}$. وخلال
الزلزال تتحرّك الأعمدة بمقدار $(s/\pi)\arctan(D/x)$: \qty{2.44}{m} عند
\qty{1}{km}، \qty{1.28}{m} عند \qty{15}{km}، \qty{0.47}{m} عند \qty{50}{km}.
(ج) المجاميع هي $0.11 + 2.44$، و$1.28 + 1.28$ و$2.08 + 0.47$، كلها
\qty{2.55}{m} $= s/2$: بعد الزلزال يصبح السور مستقيمًا من جديد على الجانب الشرقي، مزاحًا بمقدار $s/2$، والجانب الغربي، بالتماثل، بمقدار $-s/2$.
\end{solution}

\begin{solution}{exo:b1:earthquake-source:12}
(أ) $\mathbf P = (0, 1, 0)$ و$\mathbf T = (1, 0, 0)$. الزوج
$\mathbf n = (\mathbf T + \mathbf P)/\sqrt2 = (0.707, 0.707, 0)$،
$\mathbf d = (\mathbf T - \mathbf P)/\sqrt2 = (0.707, -0.707, 0)$ يحقّق
$\mathbf P = (\mathbf n - \mathbf d)/\sqrt2$ و$\mathbf T = (\mathbf n +
\mathbf d)/\sqrt2$. والمتجهتان أفقيتان، \omterm{def:b1:earthquake-source:focal-mechanism}{فالمستويان العقديان} شاقوليان: المستوى الذي عموده $\mathbf n$ (نحو $045$) يمتدّ $135$--$315$، والمستوى الذي عموده $\mathbf d$ (نحو $135$) يمتدّ $045$--$225$. (ب) على المستوى الذي يمتدّ $315$، الكتلة الواقعة على الجانب الشمالي الشرقي تتحرّك على طول $\mathbf d$، نحو $135$: فمن الكتلة الأخرى، نحو اليمين، انزلاق \omterm{def:b1:faults-folds-fabrics:fault-types}{أيمن}. وعلى المستوى الذي يمتدّ $045$، الكتلة الواقعة على الجانب الجنوبي الشرقي تتحرّك على طول $\mathbf n$، نحو $045$: انزلاق \omterm{def:b1:faults-folds-fabrics:fault-types}{أيسر}. (ج) الصدع الجديد يصنع $\ang{45} - \phi/2 =
\ang{30}$ مع $\sigma_1$: فمع $\sigma_1$ الاتجاه $000$، تمتدّ الصدوع $030$ (\omterm{def:b1:faults-folds-fabrics:fault-types}{أيسر}) و$330$ (\omterm{def:b1:faults-folds-fabrics:fault-types}{أيمن})، على \ang{15} من المستويين العقديين. \omterm{def:b1:earthquake-source:focal-mechanism}{ومحور P} ثابت على \ang{45} من المستويين العقديين بالهندسة، أيًّا كان احتكاك الصخر، فلا يتطابق مع $\sigma_1$ إلا إذا انزلق الصدع على زاوية \ang{45} منه. (د) الصدع الذي يمتدّ $040$ يبعد \ang{5} عن \omterm{def:b1:earthquake-source:focal-mechanism}{المستوى العقدي} الذي يمتدّ $045$: فهذا المستوى هو الصدع على الأرجح، والانزلاق عليه \omterm{def:b1:faults-folds-fabrics:fault-types}{أيسر}.
\end{solution}

\begin{solution}{pb:b1:earthquake-source:1}
\textbf{1.} الفواصل هي $1290 - 1050 = 240$، و$1480 - 1290 = 190$ و$1712 - 1480 = 232$ سنة؛ ومتوسطها \qty{220.7}{yr}.

\textbf{2.} متوسط الإزاحة هو $(7.0 + 6.0 + 5.5 + 7.5)/4 = \qty{6.5}{m}$، و
$T = 6.5/0.030 = \qty{216.7}{yr}$، قريبًا من متوسط الفاصل.

\textbf{3.} إزاحات الزلازل الثلاثة الأخيرة، $6.0 + 5.5 + 7.5 =
\qty{19.0}{m}$، تراكمت بين 1050 و\num{1712}، في 662 سنة: $\qty{19.0}{m}
/\qty{662}{yr} = \qty{28.7}{mm.yr^{-1}}$، أي \qty{4}{\percent} أدنى من
المعدّل الجيوديزي \qty{30}{mm.yr^{-1}}. ويتّفق الاثنان ضمن عدم يقين التواريخ
والإزاحات: فالقطعة تطلق تقريبًا كل تحميلها في زلازلها الكبيرة.

\textbf{4.} $2025 - 1712 = 313$ سنة، والعجز هو
$\qty{0.030}{m.yr^{-1}}\times\qty{313}{yr} = \qty{9.39}{m}$.

\textbf{5.} العجز يفوق كل إزاحة في الخندق، وأكبرها \qty{7.5}{m}، والزمن منذ 1712 يتجاوز أطول فاصل، 240 سنة. وبهذا المعنى تكون القطعة متأخّرة؛ ولا يعني ذلك أن الزلزال وشيك، بل فقط أن الصدع يحمل الآن انزلاقًا أكثر مما أطلقه في السجل مرة واحدة.

\textbf{6.} تتراوح الفواصل بين 190 و\num{240} سنة، والإزاحات بين \qtyrange{5.5}{7.5}{m}: فالدورة ليست دورية، ومتانة الصدع عند الزلزال التالي مجهولة. ويعطي الخندق متوسطًا وانتشارًا، ومنهما يمكن تقدير احتمال، لا تاريخ.

\textbf{7.} $A = \qty{250}{km}\times\qty{15}{km} = \qty{3750}{km^2} =
\qty{3.75e9}{m^2}$.

\textbf{8.} $M_0 = \num{3e10}\times\num{3.75e9}\times\num{9.39} =
\qty{1.06e21}{N.m}$.

\textbf{9.} $M_0 = \num{3e10}\times\num{3.75e9}\times\num{7.5} =
\qty{8.44e20}{N.m}$، إذن $M_w = \tfrac23(20.926 - 9.1) = 7.88$.

\textbf{10.} النسبة $9.39/7.5 = 1.25$، وهي $\tfrac23\log_{10}1.25 =
0.07$ من القدر.

\textbf{11.} نصف المساحة يقسم العزم على اثنين، ويخفّض القدر بمقدار
$\tfrac23\log_{10}2 = 0.20$: $M_w = 7.75$.

\textbf{12.} $\mu A V = \num{3e10}\times\num{3.75e9}\times\num{0.030} =
\qty{3.38e18}{N.m}$ في السنة.

\textbf{13.} $b = 1$. $N(5) = 10^{4.3 - 5} = 0.20$ في السنة و$N(6) = 0.020$ في
السنة.

\textbf{14.} $N(7) = 0.0020$ في السنة: واحد كل 500 سنة؛ $N(8) = 0.00020$ في
السنة: واحد كل \num{5000} سنة.

\textbf{15.} يعطي القانون زلزالًا واحدًا قدره 7.95 أو أكثر كل
$1/10^{4.3 - 7.95} = \num{4500}$ سنة، والخندق واحدًا كل 220 سنة، أي أكثر تكرارًا بعشرين مرة. والقانون مطابَق لزلازل كثيرة صغيرة من كل صدوع المنطقة؛ وإذا استُقرئ إلى الزلازل الكبيرة لقطعة واحدة، فإنه يقلّل من شأنها كثيرًا. وتكرار الزلازل الكبيرة على صدع يُؤخذ من سجلّه الخاص ومن تحميله، لا من الصغيرة.

\textbf{16.} $\num{5.5e17}/\num{3.38e18} = 0.16$: الزلازل المعتدلة
تطلق \qty{16}{\percent} من العزم الذي تتراكمه القطعة.

\textbf{17.} بحسب \cref{prop:b1:earthquake-source:moment-steps}،
$10^{1.5\times(7.95 - 6)} = 10^{2.92} \approx 840$ مرة أكثر.

\textbf{18.} مع $b = 1$، يهيمن أكبر الزلازل على عزم الزلازل: فمهما فعلت الزلازل المعتدلة، فهي تُطلق حصة صغيرة، هنا \qty{16}{\percent}، والباقي، \qty{84}{\percent}، أي عجز الانزلاق، ينتظر الزلزال الكبير.

\textbf{19.} $v(-x) = -v(x)$ لكل زوج، فالملفّ ضدّ متماثل. وبعيدًا عن الصدع تميل السرعات إلى $\pm V/2 = \pm\qty{15}{mm.yr^{-1}}$.

\textbf{20.} $D = x/\tan(\pi v/V) = 15/\tan(\pi\times7.5/30) = 15/\tan\ang{45} =
\qty{15}{km}$: عند \omterm{def:b1:earthquake-source:locking-depth}{عمق القفل} تكون السرعة $V/4$.

\textbf{21.} عند \qty{5}{km}: $D = 5/\tan(\pi\times3.1/30) = 5/\tan\ang{18.6} =
\qty{14.9}{km}$؛ وعند \qty{40}{km}: $D = 40/\tan\ang{69.6} = \qty{14.9}{km}$. والسرعات معطاة إلى \qty{0.1}{mm.yr^{-1}}، وهذا التقريب وحده يحرّك $D$ بنحو \qty{0.1}{km}؛ ومع بيانات حقيقية سيكون ضجيج كل محطة أكبر، ويُلاءَم \omterm{def:b1:earthquake-source:locking-depth}{عمق القفل} لكل المحطات معًا.

\textbf{22.} \omterm{def:b1:earthquake-source:locking-depth}{عمق القفل}، \qty{15}{km}، يساوي عمق أعمق الزلازل. فالجزء المقفل من الصدع هو جزؤه الهشّ، \omterm{def:b1:stress-strain-rheology:strength-envelope}{الطبقة المولّدة للزلازل}، حيث يمسكه الاحتكاك بين الزلازل ويفشل فيها؛ وتحته يزحف الصخر الأدفأ (\cref{ch:b1:stress-strain-rheology})، وينزلق الصدع انزلاقًا ثابتًا، وهو الانزلاق بالمعدّل $V$ تحت $D$ في النموذج.

\textbf{23.} $V/(\pi D) = \num{0.030}/(\pi\times\num{15000}) = \num{6.4e-7}$
في السنة.

\textbf{24.} مع $M_0 = \qty{1.056e21}{N.m}$، $\log_{10}M_0 = 21.024$ و
$M_w = \tfrac23(21.024 - 9.1) = 7.949$. لأقرب 0.01، فإن الزلزال الذي يطلق العجز كله في القطعة له \omterm{def:b1:earthquake-source:mw}{قدر العزم}
\textbf{7.95}.
\end{solution}
